\begin{frame}{``노이즈''가 아니라 ``시스템틱한 편차''}
  \begin{itemize}
    \item 예: 시뮬레이션의 마찰 계수가 실제보다 \textbf{3배} 작다면,
      시뮬레이션 에피소드 \emph{백만} 개를 모아도 그 오차는 \textbf{사라지지
      않는다}.
    \item Week 4(ML1)의 bias-variance 분해로 말하면:
  \end{itemize}
  \vskip0.4em
  \begin{center}
    \kb{variance}는 데이터로 줄이고, \kb{bias}는 데이터로는 줄이지 못한다.
  \end{center}
  \vskip0.6em
  \begin{itemize}
    \item bias는 ``모델을 더 잘 맞추면'' 줄일 수 있을 뿐.
    \item 그래서 시뮬레이션에서 \emph{더 많이} 학습하는 것만으로 격차가
      해결되지 않는다 --- 격차의 원인을 \textbf{구조적으로} 다루는 방법
      (무작위화, 모델 보정)이 필요.
  \end{itemize}
\end{frame}
